\documentclass[11pt]{article}
\usepackage[a4paper,margin=2cm]{geometry}
\usepackage{amsmath}
\usepackage{enumitem}
\usepackage{booktabs}
\usepackage{array}
\usepackage{longtable}
\usepackage{xcolor}
\usepackage{parskip}
\setlist{itemsep=3pt, topsep=3pt}
\renewcommand{\arraystretch}{1.3}
\newcommand{\q}[1]{\begin{quote}\itshape ``#1''\end{quote}}

\begin{document}

{\LARGE\bfseries TinyIDS: the project explained}\\[4pt]
{\large What the proposal asks for, what every term means, and how we will do it}

\vspace{6pt}
TinyIDS = a tiny \textbf{IDS (Intrusion Detection System)}: a program that watches network traffic and raises an alert when it sees an attack.
It must be small and fast enough to run on a \textbf{Raspberry Pi 4} acting as the \textbf{IoT (Internet of Things)} gateway.

\tableofcontents

%% =====================================================================
\newpage
\section{The two halves of the project: offline and live}

The proposal asks for two different kinds of work. Mixing them up causes most of the confusion.

\subsection*{(a) Offline: on the dataset}
\q{Benchmark at least four classifiers on CICIoT2023 for binary (benign vs attack), 8-class and 34-class detection.} (Objective 1)

\textbf{CICIoT2023} = \textbf{Canadian Institute for Cybersecurity IoT dataset 2023}.
Its \texttt{test.csv} already contains \textbf{all 8 families}: DDoS, DoS, Mirai, Recon, Spoofing, Web-based, BruteForce and Benign.
You do \emph{not} perform these attacks yourself: the dataset authors already ran them in their lab and recorded the traffic.
Your model is trained on \texttt{train.csv} and scored on \texttt{test.csv}. This gives the ``accuracy'' numbers of the report.

\subsection*{(b) Live: on your testbed}
\q{An attacker laptop generates SYN/UDP floods (hping3), port scans (nmap) and MQTT brute force.} (Methodology, step 5)

This is the \textbf{only} live attack list in the proposal. It covers 4 of the 8 families:

\begin{center}
\begin{tabular}{@{}l l l@{}}
\toprule
\textbf{Live attack} & \textbf{Tool} & \textbf{Dataset family} \\
\midrule
SYN flood & hping3 & DDoS / DoS \\
UDP flood & hping3 & DDoS / DoS \\
Port scan & nmap & Recon \\
MQTT brute force & small Python script & BruteForce \\
Normal sensor traffic & NodeMCUs & Benign \\
\bottomrule
\end{tabular}
\end{center}

\textbf{Mirai, Spoofing and Web-based are evaluated only on the dataset.} Nothing is missing.
(Optional extra later: ARP spoofing can be done live on your own hotspot.)

\subsection*{What the acronyms mean}
\begin{itemize}
\item \textbf{DoS} = Denial of Service; \textbf{DDoS} = Distributed Denial of Service (many attackers at once).
\item \textbf{SYN} = synchronize (the ``let's connect'' TCP packet); \textbf{TCP} = Transmission Control Protocol; \textbf{UDP} = User Datagram Protocol.
\item \textbf{hping3} = a packet generator (``ping on steroids''); \textbf{nmap} = Network Mapper, a port scanner.
\item \textbf{MQTT} = Message Queuing Telemetry Transport; \textbf{CoAP} = Constrained Application Protocol.
\end{itemize}

%% =====================================================================
\newpage
\section{Machine learning in plain English}

\subsection*{The basics}
\begin{itemize}
\item \textbf{Feature} = one input column (e.g.\ \texttt{Rate}, \texttt{syn\_flag\_number}). CICIoT2023 has 46.
\item \textbf{Label} = the correct answer for a row (e.g.\ \texttt{DDoS-SYN\_Flood}).
\item \textbf{Classifier} = a program that looks at a row's features and predicts its label. It \emph{learns} this by looking at millions of labelled rows (``training'').
\item \textbf{Binary / 8-class / 34-class} = how detailed the answer must be: attack-or-benign, which family, or which exact attack.
\end{itemize}

\subsection*{The four models the proposal asks for}
\q{Train Random Forest, XGBoost/LightGBM, a 3-layer MLP and a small 1D-CNN.} (Methodology, step 2)

\begin{longtable}{@{}>{\raggedright}p{4.2cm} >{\raggedright\arraybackslash}p{11.8cm}@{}}
\toprule
\textbf{Model} & \textbf{What it is} \\
\midrule
\endhead
\textbf{Decision tree} (building block) & A flowchart of yes/no questions learned from data:
``\texttt{Rate} $>$ 5000? $\rightarrow$ \texttt{syn\_flag\_number} $>$ 0.8? $\rightarrow$ SYN flood''. Easy to understand, but one tree alone makes mistakes. \\
\textbf{Random Forest (RF)} & \textbf{Hundreds of decision trees that vote.} Each tree is trained on a slightly different random part of the data; the majority answer wins. Accurate, robust, easy to use. \\
\textbf{XGBoost} = eXtreme Gradient Boosting & Also many trees, but built \textbf{one after another}: each new tree focuses on fixing the mistakes of the previous ones (``boosting''). Usually the most accurate on table data like ours. \\
\textbf{LightGBM} = Light Gradient Boosting Machine & Microsoft's faster, lighter version of the same boosting idea. The proposal says ``XGBoost/LightGBM'': \textbf{pick one}. We will use LightGBM (fast, small models). \\
\textbf{MLP} = Multi-Layer Perceptron & The classic \textbf{ANN (Artificial Neural Network)}: layers of ``neurons'' connected by weights. ``3-layer'' = three layers. \\
\textbf{1D-CNN} = one-dimensional Convolutional Neural Network & A neural network that slides a small window \textbf{along a row of numbers} (1-D) instead of across an image (2-D). For table data it is usually \emph{not} better than trees; it is included because many IDS papers use it, so it gives a fair comparison point. \\
\bottomrule
\end{longtable}

\textbf{Expected result:} the tree models (RF, LightGBM) will probably be the most accurate \emph{and} the fastest/smallest. That is a perfectly good finding.

\subsection*{How models are scored}
\begin{itemize}
\item \textbf{Accuracy} = fraction of rows predicted correctly. Misleading on unbalanced data (predicting ``DDoS'' for everything already gives 73\%).
\item \textbf{Precision} = of the rows the model called ``attack X'', how many really were X.
\item \textbf{Recall} = of all real ``attack X'' rows, how many the model found.
\item \textbf{F1-score} = one number combining precision and recall (their harmonic mean), from 0 to 1.
\item \textbf{Macro-F1} = the F1 of every class \textbf{averaged equally}, so rare attacks (Web, BruteForce) count as much as DDoS. The proposal uses this.
\item \textbf{Confusion matrix} = a table of ``real class vs predicted class'' showing exactly which attacks get mixed up.
\item \textbf{FPR (False Positive Rate)} = how often normal traffic is wrongly flagged as an attack.
\end{itemize}

%% =====================================================================
\newpage
\section{Making models small: compression}

The Pi is much weaker than a laptop, so the trained models are shrunk before deployment.

\begin{longtable}{@{}>{\raggedright}p{4.2cm} >{\raggedright\arraybackslash}p{11.8cm}@{}}
\toprule
\textbf{Term} & \textbf{Meaning} \\
\midrule
\endhead
\textbf{Pruning} & Cutting away the parts of a model that matter least.
Trees: fewer trees and less depth. MLP: remove the weakest connections between neurons. \\
\textbf{Quantization} & Storing the model's numbers with fewer bits, e.g.\ 32-bit decimal numbers $\rightarrow$ \textbf{INT8} (8-bit integers).
About $4\times$ smaller and faster, with a very small accuracy loss. \\
\textbf{Quantized tree ensemble} & ``Ensemble'' = many trees together (RF or LightGBM). Quantized = its thresholds/values stored in fewer bits, and usually pruned too. \\
\textbf{TFLite} = TensorFlow Lite & Google's format for running neural networks (MLP, CNN) on small devices like phones and the Pi. \\
\textbf{ONNX} = Open Neural Network Exchange & An open, universal model file format. \textbf{ONNX Runtime} runs ONNX models fast on almost any device, including the Pi. Tree models can also be exported to ONNX. \\
\bottomrule
\end{longtable}

%% =====================================================================
\section{SHAP: explaining every alert}

\textbf{SHAP} = \textbf{SHapley Additive exPlanations}. It is based on Shapley values from game theory: it treats the features like players in a team
and fairly splits ``credit'' for the decision among them.

For \textbf{one} prediction it gives every feature a score: how much it pushed the decision towards (+) or away from ($-$) the predicted class. Example alert:
\begin{verbatim}
ALERT: DDoS-SYN_Flood
  because  syn_flag_number = 1.0      (+0.45)
           Rate = 98,000 packets/s    (+0.31)
           ack_flag_number = 0.0      (+0.12)
\end{verbatim}
These are the \textbf{``top 3 features''} the dashboard must show.
\textbf{TreeSHAP} is the fast, exact version of SHAP for tree models (RF, LightGBM, XGBoost).

SHAP has \textbf{two uses} in this project:
\begin{enumerate}
\item \textbf{Explain live alerts} (RQ3, Objective 5).
\item \textbf{Rank features} (RQ1, Objective 2): average the absolute SHAP scores over many rows $\rightarrow$ a ranking of the most important features
$\rightarrow$ retrain using only the top 5 / 10 / 15 / 20.
\end{enumerate}

%% =====================================================================
\newpage
\section{The three research questions (RQ) in plain words}

\subsection*{RQ1: how many features do we really need?}
\q{How few flow features (out of 46 in CICIoT2023) are needed to keep macro-F1 within 2\% of the full-feature model?}
\begin{enumerate}
\item Train with \textbf{all 46 features}. Say macro-F1 = 0.90.
\item Retrain with only the top 5, 10, 15, 20 features (ranked by SHAP).
\item Find the \textbf{smallest} number that still reaches \textbf{at least 0.88} (within 2 points of 0.90).
\end{enumerate}
It does \textbf{not} mean ``98\% accuracy''. Fewer features = faster feature extraction on the Pi.
Result: one plot, macro-F1 vs number of features.

\subsection*{RQ2: which compressed model is best on the Pi?}
\q{Which compressed model (quantized tree ensemble, pruned MLP, small 1D-CNN) gives the best accuracy per millisecond and per megabyte on the Pi?}
Compare the three compressed models on the Pi: accuracy vs speed (milliseconds per prediction) vs size (megabytes). Result: one benchmark table.

\subsection*{RQ3: can SHAP keep up?}
\q{Can SHAP-based explanations be produced on the gateway fast enough to accompany live alerts?}
Measure how long TreeSHAP takes per alert on the Pi. Target: under 50 ms. Result: one number (plus a few examples).

\section{Targets from the proposal (Evaluation metrics)}
\begin{center}
\begin{tabular}{@{}l >{\raggedright}p{4.6cm} >{\raggedright\arraybackslash}p{7.6cm}@{}}
\toprule
\textbf{Category} & \textbf{What is measured} & \textbf{Target} \\
\midrule
Detection & macro-F1, per-class recall, FPR & \textbf{Binary F1 $\geq$ 0.98}; 8-class within 2\% of full-feature model \\
Efficiency & model size, peak RAM on the Pi & \textbf{$<$ 1 MB} model; \textbf{$<$ 100 MB} RAM \\
Speed & latency per flow, throughput & \textbf{$<$ 5 ms} per flow; \textbf{$\geq$ 500 flows/s} \\
Energy & mJ (millijoules) per inference & just report it for every model \\
Explainability & extra SHAP time per alert & \textbf{$<$ 50 ms} \\
\bottomrule
\end{tabular}
\end{center}
Note: RAM is in \textbf{megabytes}, not kilobytes. \textbf{RAM} = Random Access Memory. For a Pi 4 with 4 GB RAM these targets are relaxed;
tree models easily fit in under 1 MB.

%% =====================================================================
\newpage
\section{The testbed: build everything on the laptop first}

\subsection*{The plan}
\begin{center}
\begin{tabular}{@{}l >{\raggedright}p{6.8cm} l l@{}}
\toprule
\textbf{Phase} & \textbf{What} & \textbf{Where} & \textbf{Needs lab/Pi?} \\
\midrule
1. Train & Train RF, LightGBM, MLP, 1D-CNN on CICIoT2023 (binary + 8-class), SHAP ranking, retrain on top-k features, save model files. & Laptop / Colab & No \\
2. Laptop testbed & Fake gateway + fake sensors + attacker as Docker containers; run attacks; check detection. & Laptop (Docker) & No \\
3. Pi & Copy the same code to the Pi, real NodeMCUs, measure size / RAM / speed / energy. & Lab & \textbf{Yes} \\
4. Finish & Dashboard with top-3 SHAP features, demo video, report. & Laptop & No \\
\bottomrule
\end{tabular}
\end{center}
Only phase 3 needs the lab. The Pi runs Linux and Python like the Docker containers, so the code moves over almost unchanged.
(\textbf{Colab} = Google Colaboratory, free online Python notebooks with a GPU, Graphics Processing Unit.)

\subsection*{The laptop testbed (phase 2)}
\textbf{Docker} runs small isolated Linux ``machines'' (\textbf{containers}) on the laptop, joined by a \textbf{private virtual network}:

\begin{verbatim}
        Docker private network (exists only inside the laptop)
 +--------------+   +----------------------------------+   +--------------+
 | sensors x3   |-->| GATEWAY  (stands in for the Pi)  |<--| ATTACKER     |
 | fake NodeMCU |   | Mosquitto + sniffer + IDS model  |   | hping3, nmap |
 | publishers   |   +----------------------------------+   | brute force  |
 +--------------+                                           +--------------+
\end{verbatim}
\begin{itemize}
\item \textbf{Gateway}: runs Mosquitto (the MQTT broker), captures every packet, builds the features (Rate, SYN fraction, \dots) for each window of packets, and asks the model ``attack or not?''. Alerts go to the dashboard with the top-3 SHAP features.
\item \textbf{Sensors}: small Python scripts publishing fake temperatures over MQTT, like the NodeMCU test.
\item \textbf{Attacker}: has hping3, nmap and an MQTT brute-force script.
\end{itemize}
\textbf{Big advantage:} attack traffic never leaves the laptop, so it is safe to run anywhere (hostel, home) without touching the college network.
Already available on the laptop: Docker Desktop, \textbf{WSL} (Windows Subsystem for Linux) with Ubuntu, Wireshark/Npcap, 15 GB RAM.

\subsection*{What can be done without the lab or the Pi}
\begin{itemize}
\item All of phase 1 (training, SHAP, feature selection, compression, accuracy tables and plots).
\item All of phase 2 (live detection on the Docker testbed, all four live attacks).
\item The dashboard and most of the report.
\item The NodeMCUs can join the laptop testbed over a phone hotspot (laptop runs the broker and IDS) whenever you have them with you.
\end{itemize}
Only the \textbf{Pi-specific measurements} (model speed, RAM and energy on the Pi) need the Pi.

%% =====================================================================
\newpage
\section{Performing the attacks}

\subsection*{No device has to be ``compromised''}
The attacker simply sends bad traffic \textbf{at} the gateway, the way a real attacker would.
The gateway does not have to be insecure: a flood still hits it, and the IDS's job is to \textbf{notice}.

\begin{longtable}{@{}>{\raggedright}p{3.2cm} >{\raggedright}p{3.8cm} >{\raggedright\arraybackslash}p{8.8cm}@{}}
\toprule
\textbf{Attack} & \textbf{Tool / command idea} & \textbf{What happens} \\
\midrule
\endhead
SYN flood & \texttt{hping3 -S --flood -p 1883 <gateway>} & Thousands of ``let's connect'' (SYN) packets per second that never complete the handshake. \\
UDP flood & \texttt{hping3 --udp --flood <gateway>} & A blast of junk UDP packets. \\
Port scan & \texttt{nmap <gateway>} & Knocks on every port to see which are open (closed ports reply with RST). \\
MQTT brute force & Python script & Tries many username/password pairs on Mosquitto (needs a password set on the broker). \\
\bottomrule
\end{longtable}

\subsection*{What the professor probably meant by ``compromise a device''}
Real botnets like Mirai first take over IoT devices and then use them to flood a target. You do not have to reproduce that:
the dataset already contains Mirai traffic. But you \emph{can} show the idea safely with your own NodeMCUs (next section).

\subsection*{Can we install real Mirai on the NodeMCU?}
\textbf{No.} Mirai is Linux malware for cameras and routers that run a small Linux system. The \textbf{ESP8266} chip on the NodeMCU
does not run Linux, so Mirai cannot run on it. Real Mirai would also automatically try to break into other devices on the network,
which on the college network would be a genuine attack on other people's devices. We imitate its behaviour instead.

\subsection*{Rules}
\begin{itemize}
\item Attack \textbf{only your own setup}: the Docker network on the laptop, or the Pi + NodeMCUs on your own phone hotspot.
\item \textbf{Never} run floods or scans on the college network: IT will see a real attack.
\end{itemize}

%% =====================================================================
\newpage
\section{Planned: laptop as ``master'' of the NodeMCUs (simulated botnet)}

\subsection*{The idea}
This imitates how Mirai controls infected devices, using only your own hardware:
\begin{enumerate}
\item Each NodeMCU runs a sketch that behaves like a \textbf{normal sensor}: it publishes data every few seconds.
\item It also \textbf{listens} on an MQTT command topic, e.g.\ \texttt{cmd/all}.
\item From the laptop (the ``master'', like Mirai's \textbf{C2} = Command-and-Control server) you publish a command, e.g.\\
\texttt{\{"attack":"udp\_flood","seconds":10\}}
\item Every NodeMCU suddenly starts attacking the gateway at the same time $\Rightarrow$ a small \textbf{DDoS} from ``compromised sensors''.
\item After the time is up (or on a \texttt{stop} command) they go back to being normal sensors.
\item The IDS should flag the flood and show which features gave it away.
\end{enumerate}
\textbf{Yes, this is possible} and it makes a strong demo: normal traffic $\rightarrow$ command $\rightarrow$ coordinated attack $\rightarrow$ alert.

\subsection*{What a NodeMCU can and cannot do}
\begin{center}
\begin{tabular}{@{}>{\raggedright}p{4.6cm} >{\raggedright}p{3.8cm} >{\raggedright\arraybackslash}p{7.2cm}@{}}
\toprule
\textbf{Attack from a NodeMCU} & \textbf{Possible?} & \textbf{Note} \\
\midrule
UDP flood & Yes & Like Mirai's \texttt{udpplain}. A few thousand packets/s per board. \\
TCP connection flood & Yes & Opens many real connections quickly. \\
Raw SYN flood with fake IPs & No & The ESP8266 cannot build raw packets; use hping3 on the laptop for this. \\
HTTP (web request) flood & Yes, if the gateway runs a web page & Optional. \\
\bottomrule
\end{tabular}
\end{center}

\subsection*{Built-in safety (part of the plan)}
\begin{itemize}
\item The attack target is \textbf{fixed in the sketch} (only the gateway's IP), not taken from the command.
\item Every attack has a \textbf{maximum duration} (e.g.\ 30 s) and stops by itself.
\item Use it only on your own hotspot / testbed network.
\end{itemize}

\subsection*{Where it fits}
After phase 2 works with hping3, we add this as the ``compromised IoT device'' demo in phase 3 (or on the laptop testbed with NodeMCUs on a hotspot).

%% =====================================================================
\section{Next steps}
\begin{enumerate}
\item \textbf{Training script}: load a sample of \texttt{train.csv}, keep features we can compute live, train RF + LightGBM + MLP + 1D-CNN for binary and 8-class, report macro-F1 and confusion matrices on \texttt{test.csv}, save the models.
\item SHAP feature ranking and retraining with top 5/10/15/20 (RQ1 plot).
\item Docker testbed with gateway, sensors and attacker.
\item Live detection + SHAP alerts + dashboard.
\item Pi deployment and benchmark table (RQ2, RQ3).
\item Optional: master-controlled NodeMCU botnet demo.
\end{enumerate}

%% =====================================================================
\newpage
\section*{Appendix: full forms of all abbreviations}
\addcontentsline{toc}{section}{Appendix: full forms of all abbreviations}
\begin{longtable}{@{}>{\bfseries}l >{\raggedright\arraybackslash}p{12.5cm}@{}}
\toprule
Abbreviation & Full form / meaning \\
\midrule
\endhead
ACK & Acknowledge (TCP flag: ``got it'') \\
ANN & Artificial Neural Network \\
ARP & Address Resolution Protocol (``who has this IP address?'') \\
C2 & Command-and-Control (the server that gives orders to a botnet) \\
CICIoT2023 & Canadian Institute for Cybersecurity IoT dataset 2023 \\
CNN / 1D-CNN & Convolutional Neural Network / one-dimensional CNN \\
CoAP & Constrained Application Protocol \\
DDoS & Distributed Denial of Service \\
DNS & Domain Name System (turns names like google.com into IP addresses) \\
DoS & Denial of Service \\
ESP8266 & The Wi-Fi microcontroller chip on the NodeMCU (made by Espressif) \\
F1 & F1-score (harmonic mean of precision and recall) \\
FPR & False Positive Rate \\
GPU & Graphics Processing Unit \\
GRE & Generic Routing Encapsulation (a tunnelling protocol, used by Mirai floods) \\
IAT & Inter-Arrival Time (gap between two packets) \\
ICMP & Internet Control Message Protocol (ping) \\
IDS & Intrusion Detection System \\
INT8 & 8-bit integer \\
IoT & Internet of Things \\
IP & Internet Protocol \\
JSON & JavaScript Object Notation (text format for data) \\
LightGBM & Light Gradient Boosting Machine \\
MITM & Man-in-the-Middle \\
MLP & Multi-Layer Perceptron \\
MQTT & Message Queuing Telemetry Transport \\
NodeMCU & Node MicroController Unit (the ESP8266 development board) \\
ONNX & Open Neural Network Exchange \\
RAM & Random Access Memory \\
RF & Random Forest \\
RQ & Research Question \\
RSSI & Received Signal Strength Indicator \\
RST & Reset (TCP flag: ``abort this connection'') \\
SHAP & SHapley Additive exPlanations \\
SQL & Structured Query Language (databases) \\
SYN & Synchronize (TCP flag: ``let's connect'') \\
TCP & Transmission Control Protocol \\
TFLite & TensorFlow Lite \\
TTL & Time To Live \\
UDP & User Datagram Protocol \\
WSL & Windows Subsystem for Linux \\
XGBoost & eXtreme Gradient Boosting \\
XSS & Cross-Site Scripting \\
\bottomrule
\end{longtable}

\end{document}
